%% Copyright (c) 2018-2026 Weitian LI <wt@liwt.net>
%% CC BY 4.0 License
%%
%% Revised for Yang Desheng, 2026

\documentclass[zh]{resume}

\fileinfo{
  \faEdit{} \today
}

\name{德生}{杨}

\profile{
  \mobile{187-7975-0679}
  \email{desyang@qq.com}
  \gitee{desyang-hub} \\
  \university{南昌航空大学}
  \degree{计算机科学与技术 \textbullet 硕士（在读）}
  \birthday{2000-09}
  \address{江西 南昌}
}

\begin{document}
\makeheader

%======================================================================
% 摘要（Summary）
%======================================================================
\begin{abstract}
扎实掌握现代 C++（C++11/14）核心特性（智能指针、模板、多线程、RAII）及 Linux 系统编程，具备高性能、高可靠系统软件开发能力；
熟练使用 GDB/Valgrind/perf 等工具进行性能分析与内存调试，熟悉 CMake 构建体系与 Git 工作流；
LeetCode 刷题 300+，获省级 ACM 一等奖，算法与工程实现能力兼备；研究方向聚焦文档图像智能处理，已授权发明专利 3 项。
\end{abstract}

%======================================================================
\sectionTitle{求职意向}{\faBullseye}
%======================================================================
\begin{itemize}
  \item C/C++ 系统开发工程师（Linux / 后端 / 基础软件方向）
\end{itemize}

%======================================================================
\sectionTitle{技术能力}{\faWrench}
%======================================================================
\begin{competences}
  \comptence{编程语言}{
    C++ (C++11/14, STL, 多线程, 智能指针), 
    Python (NumPy, Pandas, 脚本自动化)
  }
  \comptence{系统与工具}{
    Linux 系统编程, GDB/Valgrind/perf, CMake/Make, Git, SSH, Docker（基础）
  }
  \comptence{算法与数据结构}{
    熟练掌握常见数据结构与算法，LeetCode 300+，ACM 省赛一等奖
  }
  \comptence{深度学习}{
    PyTorch, OpenCV, 图像预处理/后处理，轻量化模型部署
  }
  \comptence{\icon{\faLanguage} 语言}{
    英语 CET-4，可流畅阅读英文技术文档
  }
\end{competences}

%======================================================================
\sectionTitle{教育背景}{\faGraduationCap}
%======================================================================
\begin{educations}
  \education%
    {2024.09}%
    [至今]%
    {南昌航空大学}%
    {信息工程学院}%
    {计算机科学与技术}%
    {硕士（在读）}

  \separator{0.5ex}
  \education%
    {2020.09}%
    [2024.06]%
    {南昌航空大学}%
    {信息工程学院}%
    {计算机科学与技术}%
    {工学学士}
\end{educations}

%======================================================================
\sectionTitle{项目经历}{\faCode}
%======================================================================
\begin{itemize}
  \item \textbf{高性能异步日志库（C++）} \hfill 2025.03 – 2025.05 \\
        - 基于环形缓冲区与双缓冲机制设计异步写入架构，吞吐量达 50万行/秒，CPU 占用降低 35% \\
        - 支持日志分级、文件轮转、压缩归档，内存峰值控制在 10MB 以内 \\
        - 使用 CMake 构建，提供头文件即用接口，已在实验室 3 个 C++ 项目中集成

  \item \textbf{文档图像表格结构识别系统}（硕士课题）\hfill 2024.09 – 至今 \\
        - 构建轻量级 CNN 模型（PyTorch）用于表格类型判定（有线/无线/混合），准确率达 92.3% \\
        - 设计规则引擎对模型输出进行后处理，显著提升复杂表格结构还原精度 \\
        - 成果支撑 3 项发明专利（见科研成果），代码已开源至 Gitee
\end{itemize}

%======================================================================
\sectionTitle{实习经历}{\faBriefcase}
%======================================================================
\begin{experiences}
  \experience%
    [2023.05]%
    {2023.08}%
    {IT 开发工程师 @ 南昌华勤电子科技有限公司}%
    [\begin{itemize}
      \item 使用 Python (Pandas + Matplotlib) 开发自动化报表系统，日均处理 10k+ 条生产数据，人工分析时间减少 80%
      \item 配置 GitLab CI 流水线，实现代码合并后自动构建、测试并部署至 Docker 测试环境
      \item 基于 Flask 开发 RESTful 接口，支持设备状态实时查询，平均响应时间 < 150ms
    \end{itemize}]
\end{experiences}

%======================================================================
\sectionTitle{科研成果}{\faAtom}
%======================================================================
\begin{itemize}
  \setlength{\itemsep}{0.4ex}
  \renewcommand{\labelitemi}{\color{gray}\tiny$\bullet$}
  
  \item \textbf{发明专利：一种矩形印章文字识别方法}（第一发明人）\hfill 2024（已授权）\\
        专利号：ZL 2024 1 0373918.4 
  \item \textbf{发明专利：一种文档图像表格类型判定方法}（第二发明人）\hfill 2025（已授权）\\
        专利号：ZL 2025 1 0865310.8 
  \item \textbf{发明专利：一种轻量级表格结构识别方法}（第三发明人）\hfill 2025（已授权）\\
        专利号：ZL 2025 1 0409433.0     
\end{itemize}

%======================================================================
\sectionTitle{荣誉奖项}{\faTrophy}
%======================================================================
\begin{itemize}
  \item \textbf{江西省大学生程序设计竞赛 一等奖} \hfill 2022
  \item \textbf{睿抗机器人开发者大赛（RAICOM）CAIP 编程设计赛 全国二等奖} \hfill 2023
  \item 国家励志奖学金 \hfill 2021
  \item 优秀研究生 \hfill 2025
\end{itemize}

\end{document}